\pdfmapfile{+symbols.map}
\pdfmapfile{+cm.map}
\pdfmapfile{+cmextra.map}
\pdfmapfile{+lm.map}
\documentclass[11pt,letterpaper]{article}
\usepackage[margin=1in]{geometry}
\usepackage{amsmath,amssymb,amsthm}
\usepackage[T1]{fontenc}
\usepackage{lmodern}
\usepackage{microtype}
\usepackage[colorlinks=true,linkcolor=blue,urlcolor=blue,citecolor=blue]{hyperref}
\newtheorem{theorem}{Theorem}
\newtheorem{lemma}{Lemma}
\newcommand{\Q}{\mathbb{Q}}
\newcommand{\crk}{\operatorname{crk}}
\newcommand{\ncrk}{\operatorname{ncrk}}
\newcommand{\Mat}{\operatorname{Mat}}
\title{Disproof of Conjecture 00000005207}
\author{}
\date{20 October 2026}
\begin{document}
\maketitle
\vspace{-1.5em}
\begin{abstract}
The difference between noncommutative and commutative rank is not
multiplicative under tensor products. Over $\Q$, the space of skew-symmetric
$3\times3$ matrices has ranks $3$ and $2$, respectively. Tensoring with the
$1\times1$ scalar matrix space leaves this gap equal to $1$, whereas the
product of the two factor gaps is $0$. We give explicit rank certificates
and a Lean 4 proof using only the standard library.
\end{abstract}

\section{The clause being disproved}
The posted conjecture\footnote{\href{https://github.com/The-Last-Math-Competition/The-Last-Math-Competition/blob/main/conjectures/00000005207.md}{Official conjecture 00000005207 on GitHub}, checked 20 October 2026.}
states, in particular, that ``the gap is multiplicative under tensor
products.'' With
\[
 \Delta(\mathcal X)=\ncrk(\mathcal X)-\crk(\mathcal X),
\]
this clause asserts
\begin{equation}\label{eq:claim}
 \Delta(\mathcal X\otimes\mathcal Y)
 =\Delta(\mathcal X)\Delta(\mathcal Y).
\end{equation}
Both factors and their tensor product in our counterexample have matrix
size at most three. Thus the counterexample remains within the stated
low-order range. It disproves the conjecture as a conjunction; it does
not dispute the separate assertion that gaps $0$ and $1$ occur at these
sizes. No interpretation of the additional phrase about a ``multiplicative
deficit'' is needed to refute the explicit equality \eqref{eq:claim}.

\section{Rank conventions}
For a rational matrix space $\mathcal X\subseteq\Mat_n(\Q)$, use
\begin{align*}
 \crk(\mathcal X)&=\max_{M\in\mathcal X}\operatorname{rank}(M),\\
 \mathcal X^{\{d\}}&=\mathcal X\otimes\Mat_d(\Q),\\
 \ncrk(\mathcal X)&=\max_{d\geq1}
      \frac{\max_{M\in\mathcal X^{\{d\}}}\operatorname{rank}(M)}{d}.
\end{align*}
These are standard matrix-space definitions. Since $\Q$ is infinite,
commutative rank agrees with the rank of a generic linear matrix over
its rational-function field. The normalized blowup definition agrees
with free-skew-field rank; see Derksen--Makam, Definition 1.3,
Lemma 1.4, Definition 1.5, Lemma 1.8 and Definition 1.9\cite{DM}.

Concretely, if $X_0,\ldots,X_{m-1}$ generate $\mathcal X$, its elements
are $\sum_i a_iX_i$ and its $d$-blowup elements are
$\sum_i X_i\otimes T_i$, with $a_i\in\Q$ and $T_i\in\Mat_d(\Q)$.
We will prove both a universal upper bound and an attaining witness
for each rank, so no unproved existence of a maximizing witness is
used in the calculation.

\section{Explicit rank certificates}
Let
\[
 \mathcal S=\left\{
 S(a,b,c)=\begin{pmatrix}
 0&a&b\\-a&0&c\\-b&-c&0
 \end{pmatrix}:a,b,c\in\Q\right\},
 \qquad \mathcal U=\Mat_1(\Q).
\]
Write $X_0=S(1,0,0)$, $X_1=S(0,1,0)$ and $X_2=S(0,0,1)$.

\begin{lemma}\label{lem:comm}
$\crk(\mathcal S)=2$.
\end{lemma}
\begin{proof}
Direct expansion gives
$\det S(a,b,c)=-abc+abc=0$ for every $a,b,c\in\Q$.
Therefore every matrix in $\mathcal S$ has rank at most two.
The leading $2\times2$ minor of $S(1,0,0)$ has determinant $1$,
so rank two is attained.
\end{proof}

\begin{lemma}\label{lem:noncomm}
$\ncrk(\mathcal S)=3$.
\end{lemma}
\begin{proof}
Every element of $\mathcal S^{\{d\}}$ is $3d\times3d$, so its rank
is at most $3d$. This proves the upper bound $3$ simultaneously for
all positive integers $d$.
For attainment at $d=2$, take
\[
 T_0=\begin{pmatrix}1&0\\0&0\end{pmatrix},\quad
 T_1=\begin{pmatrix}0&1\\1&0\end{pmatrix},\quad
 T_2=\begin{pmatrix}0&0\\0&1\end{pmatrix}.
\]
Their block substitution is
\begin{equation}\label{eq:blowup}
 B=\sum_{i=0}^2 X_i\otimes T_i
 =\begin{pmatrix}
 0&0&1&0&0&1\\
 0&0&0&0&1&0\\
 -1&0&0&0&0&0\\
 0&0&0&0&0&1\\
 0&-1&0&0&0&0\\
 -1&0&0&-1&0&0
 \end{pmatrix}.
\end{equation}
It is nonsingular. Indeed, if $Bv=0$, rows $2,3,4,5$ give
$v_5=v_1=v_6=v_2=0$; rows $1,6$ then give $v_3=v_4=0$.
Equivalently, direct determinant expansion yields $\det B=1$,
which is the exact certificate checked in Lean.
Thus $\operatorname{rank}(B)/2=6/2=3$, proving attainment.
\end{proof}

\begin{lemma}\label{lem:unit}
$\crk(\mathcal U)=\ncrk(\mathcal U)=1$ and
$\mathcal S\otimes\mathcal U=\mathcal S$ under the usual scalar
Kronecker identification.
\end{lemma}
\begin{proof}
For commutative rank, the size bound is one and $[1]$ attains it.
For noncommutative rank, every $d$-blowup matrix has size $d$, while
$d=1$ and $[1]$ already attain normalized rank one.
Finally $X_i\otimes[1]=X_i$ entrywise for each $i$, so the coefficient
matrices of the tensor product are precisely the same three matrices.
This is equality of the actual coefficient tensors after the standard
index identification, not an assumed invariance of an unspecified rank.
\end{proof}

\begin{theorem}
Conjecture 00000005207 is false.
\end{theorem}
\begin{proof}
The preceding lemmas give
\[
 \Delta(\mathcal S\otimes\mathcal U)=3-2=1,
 \qquad
 \Delta(\mathcal S)\Delta(\mathcal U)=(3-2)(1-1)=0.
\]
Since $1\ne0$, equation \eqref{eq:claim} fails.
\end{proof}

\section{Formal coverage and reproducibility}
The file \texttt{RankGap.lean} imports only \texttt{Std} and checks
with Lean 4.31.0. Its determinant is ordinary recursive Laplace
expansion. Matrix rank bounds quantify over every minor selected by
strictly increasing row and column index lists. The predicates
\texttt{CommutativeRank} and \texttt{NoncommutativeRank} express exact
rank through an upper bound and an attaining nonsingular minor.
In particular, the noncommutative upper bound quantifies over \emph{all}
positive blowup sizes and all rational matrix substitutions.

The formal development verifies the symbolic determinant identity,
the $2\times2$ minor, the $6\times6$ determinant, the two scalar-space
ranks, and the coefficient tensor-unit equality. Its theorem
\texttt{conjecture\_00000005207\_false} negates
\texttt{MultiplicativeGap}, which requires the genuine rank predicates
for both factors and their tensor product. It restricts even the
product matrix size to at most three.

The formalization uses the standard equivalent tensor-blowup definition
of noncommutative rank directly. It does not construct free skew fields
or reprove the general equivalence theorem between these definitions;
that mathematical correspondence is cited above, not inserted as a Lean
axiom. All concrete rank calculations and the contradiction are checked
without assumptions about their values. The kernel axiom dependency
list is exactly \texttt{propext}, \texttt{Classical.choice}, and
\texttt{Quot.sound}. There are no custom axioms, admitted proofs, or
native decision shortcuts.

Run Lean with \texttt{-DwarningAsError=true} on \texttt{RankGap.lean}, or use
\texttt{lake build} with the supplied toolchain file. The accompanying
README records the theorem map and checking commands.

\begin{thebibliography}{1}
\bibitem{DM}
Harm Derksen and Visu Makam.
\emph{On non-commutative rank and tensor rank}.
\textit{Linear and Multilinear Algebra} 66 (2018), 1069--1084.
\href{https://arxiv.org/abs/1606.06701}{arXiv:1606.06701},
Definitions 1.3, 1.5 and 1.9; Lemmas 1.4 and 1.8.
The skew-symmetric rank-gap example is classical; the present disproof
uses its elementary tensor-unit obstruction.
\end{thebibliography}
\end{document}
